\begin{table*}[!htb]
\centering
\caption{Estimativas do expoente de $P(k) \sim C k^{-\gamma}$ no subgrafo e no Orkut completo.}
\label{tab:powerlaw}
\begin{adjustbox}{max width=\linewidth}
\begin{tabular}{lcc}
\toprule
Método & Subgrafo G27 & Orkut completo \\
\midrule
Regressão log-log, pontos brutos: γ (R²) & 1,61 (0,74) & 2,13 (0,84) \\
Regressão log-log, bins logarítmicos: γ (R²) & 0,94 (0,57) & 1,92 (0,89) \\
Regressão na CCDF: γ = 1 − inclinação (R²) & 2,85 (0,74) & 2,93 (0,97) \\
MLE discreta: γ ± erro-padrão & 6,00 ± 0,31 & 2,96 ± 0,004 \\
k\_min (mínima distância KS) / vértices na cauda & 119 / 259 & 168 / 289.226 \\
Distância KS na cauda & 0,057 & 0,016 \\
Vuong vs. exponencial: z (p) & -2,91 (0,0036) & 33,9 (< 1e-200) \\
Modelo favorecido na cauda & exponencial & lei de potência \\
KS global da Poisson (λ = grau médio) & 0,44 & 0,58 \\
\bottomrule
\end{tabular}
\end{adjustbox}
\end{table*}
